% =========================================================================
% introduction.tex
% =========================================================================

\section{Introduction}
\label{sec:introduction}

The Key-Value (KV) cache is the dominant memory bottleneck during
autoregressive inference in transformer-based language
models~\cite{radford2019gpt2}.
As sequence lengths grow, the cache scales linearly in both the number of
layers and the context window, consuming memory that often exceeds the
model weights themselves.
Several recent works address this bottleneck through vector quantization
of KV embeddings.

TurboQuant~\cite{zandieh2025turboquant} applies a random orthogonal
rotation, implemented via a Walsh-Hadamard Transform (WHT) or QR
decomposition, to KV vectors, inducing a concentrated
$\mathrm{Beta}(0.5,\, d/2 - 0.5)$ distribution on each coordinate.
Pre-computed Lloyd-Max scalar quantizers~\cite{lloyd1982least} then
compress each coordinate independently, achieving near-optimal
distortion within a factor of ${\approx}2.7$ of the Shannon
bound~\cite{shannon1959coding}.
A secondary one-bit Quantized Johnson-Lindenstrauss
(QJL)~\cite{zandieh2024qjl} residual correction removes the bias
that Mean Squared Error (MSE)-optimal quantizers introduce in inner-product
estimation.
PolarQuant~\cite{han2025polarquant} extends this line by transforming
embeddings into polar coordinates whose angles follow a tightly bounded
distribution, eliminating explicit normalization overhead.
Memory Caching~\cite{behrouz2026mc} proposes a complementary approach:
rather than compressing vectors, it augments recurrent models with
checkpointed hidden-state caches, interpolating between the fixed
memory of Recurrent Neural Networks (RNNs) and the growing memory of
transformers.

\paragraph{The gap.}
All four papers validate their methods exclusively on models with
attention head dimension $\dhead \geq 128$: Llama-3.1-8B, Mistral-7B,
and Qwen-2.5-7B for TurboQuant and PolarQuant; models of 760M to 1.3B
parameters trained on 30--100 billion tokens for Memory Caching.
No published evaluation exists at $\dhead = 64$, which is the head
dimension of GPT-2~\cite{radford2019gpt2} and several other
widely-studied architectures.
This omission matters because TurboQuant's theoretical guarantees rely on
Central Limit Theorem (CLT) convergence of the rotated coordinates,
and this convergence weakens as dimensionality decreases.
Community implementations have informally observed degraded performance
at $\dhead = 64$~\cite{amesianx2025turboquant, onlyterp2025turboquant},
but no systematic study has characterized the failure modes, proposed
principled corrections, or tested cross-architecture combinations at
this scale.

\paragraph{Proxy evaluation scope.}
Evaluating KV cache quantization on GPT-2 Medium exposes an inherent methodological proxy limitation. KV cache quantization is designed to alleviate memory pressure during massive, long-context retrieval, whereas GPT-2 Medium is hard-capped at a context window of only 1{,}024 tokens. Consequently, this study does not present a direct solution for large-scale production deployment bottlenecks. Instead, this work is framed as an isolated investigation into the geometric and statistical failure modes of vector quantization when compressed into small, low-dimensional spaces ($\dhead = 64$).

\paragraph{Contributions.}
This paper presents six contributions, all evaluated on GPT-2
Medium (345M parameters, $\dhead = 64$) using a single NVIDIA RTX~4060
Laptop GPU with PyTorch~2.5.1:

\begin{enumerate}[noitemsep,topsep=1pt,parsep=0pt,partopsep=0pt]

\item \textbf{Adaptive QJL Fallback Heuristic.}
I implement a simple empirical fallback heuristic that disables the QJL residual correction when the head dimension is below $\dhead = 128$. Because QJL's variance doubles at $\dhead = 64$, bypassing the residual calculation entirely for low dimensions is a practical engineering solution that prevents performance degradation. At four-bit precision with $3.2{\times}$ compression, this fallback raises nearest-neighbor Recall@1 from 38.8\% (flat TurboQuant) to 67.4\% (WHT + Asymmetric Lloyd-Max).
% Source: results/nn\_search.json → flat\_turboquant/4/recall\_1 = 38.8;
%         results/nn\_search.json → wht\_asym/4/recall\_1 = 67.4

\item \textbf{Outlier Channel Splitting at 2.5-bit.}
I show that allocating extra precision to outlier channels yields a
Pareto improvement over flat three-bit quantization: lower short-context
average perplexity (4.68 versus 4.69) and a smaller cache footprint
(957\,KB versus 1{,}218\,KB).
% Source: results/outlier\_results.json → average\_perplexity
%         (outlier\_2.5 = 4.676, flat\_3 = 4.691);
%         average\_cache\_size\_kb (outlier\_2.5 = 957, flat\_3 = 1218)

\item \textbf{FibQuant negative result.}
A joint two-dimensional radial-angular codebook in the style of
PolarQuant/FibQuant degrades to only four angular directions at four-bit
under $d = 64$, and scalar MSE quantization outperforms it at every
tested bit budget.
% Source: README.md → Finding 5; results/fibquant\_results.json

\item \textbf{MC-TurboQuant.}
I combine WHT + Asymmetric vector quantization with Gated Residual
Memory (GRM) from the Memory Caching framework, producing the first
(to my knowledge) cross-architecture fusion of these two systems.
The combined scheme maintains 97.5\% average Needle-in-a-Haystack
(NIAH) recall across context lengths 256--1{,}024 with $3.2{\times}$
compression and no perplexity regression relative to the non-MC
quantized variant.
% Source: results/mc\_niah.json → mc\_turboquant average;
%         turboquant/validate\_paper.py Test 7 output (97.50\%)

\item \textbf{Rotation profiling.}
I profile WHT and QR rotation latencies at $d = 64$ and show that
unfused WHT in pure PyTorch is slower than naive quantization, despite
the theoretical latency advantage of fast Hadamard transforms. This finding
demonstrates the necessity of fused CUDA kernels for practical
deployment.
% Source: results/rotation\_profile.json; README.md → Limitation bullet 2

\item \textbf{Checkpoint versus independent compressor ablation.}
I compare the checkpointed and independent segment-state variants of
Memory Caching, finding that independent compression achieves 100\%
NIAH recall at most positions with lower perplexity, while the
checkpoint variant suffers state drift that degrades both metrics.
% Source: results/mc\_compressor\_ablation.json

\end{enumerate}

\noindent
All experiments, data, and code are publicly available at \url{https://github.com/G26karthik/kvquant-lab}.
Section~\ref{sec:background} reviews the four source papers.
Section~\ref{sec:related} positions this work relative to the broader
KV-cache quantization literature and informal community findings.
Section~\ref{sec:method} details each of the six contributions.
Section~\ref{sec:validation} summarizes the nine theorem-validation tests.
Sections~\ref{sec:setup} and~\ref{sec:results} describe the experimental
setup and results.
Section~\ref{sec:limitations} states the limitations of this study,
and Section~\ref{sec:conclusion} concludes.
